\newpage
\phantomsection
\label{prf:rational-square-bracketing}
\LRAProofFor{thm:rational-square-bracketing}

\begin{remark*}[Return]
\hyperref[thm:rational-square-bracketing]{Return to Theorem}
\end{remark*}

\begin{theorem*}[Rational Square Bracketing]
For every $r\in\mathbb{Q}$ with $r>0$, there exist $a,b\in\mathbb{Q}$ with
$a>0$ and $b>0$ such that
\[
a^2<r<b^2.
\]
\end{theorem*}

\begin{proof}[Professional Standard Proof]
\LRAProofBodyStart
Let
\[
a:=\frac{r}{r+1}
\qquad\text{and}\qquad
b:=r+1.
\]
Since $r\in\mathbb{Q}$ and $r>0$, both $a$ and $b$ are rational and positive.

We first show that $a^2<r$. Because $r>0$, we have $r+1>1$, hence
\[
0<a=\frac{r}{r+1}<r.
\]
Also $a<1$, since
\[
\frac{r}{r+1}<1.
\]
Multiplying the inequality $a<1$ by the positive number $a$ gives
\[
a^2<a.
\]
Since $a<r$, it follows that
\[
a^2<r.
\]

Now we show that $r<b^2$. Since $b=r+1$ and $r>0$, we have
\[
b=r+1>r.
\]
Because $b>1>0$, multiplying the inequality $b>1$ by the positive number $b$
gives
\[
b^2>b.
\]
Combining this with $b>r$, we obtain
\[
r<b^2.
\]

Therefore there exist positive rational numbers $a$ and $b$ such that
\[
a^2<r<b^2,
\]
as required.
\end{proof}

\begin{proof}[Detailed Learning Proof]
\LRAProofBodyStart
We expand the professional proof by following the same sequence of justified moves.

Let
\[
a:=\frac{r}{r+1}
\qquad\text{and}\qquad
b:=r+1.
\]
Since $r\in\mathbb{Q}$ and $r>0$, both $a$ and $b$ are rational and positive.

We first show that $a^2<r$. Because $r>0$, we have $r+1>1$, hence
\[
0<a=\frac{r}{r+1}<r.
\]
Also $a<1$, since
\[
\frac{r}{r+1}<1.
\]
Multiplying the inequality $a<1$ by the positive number $a$ gives
\[
a^2<a.
\]
Since $a<r$, it follows that
\[
a^2<r.
\]

Now we show that $r<b^2$. Since $b=r+1$ and $r>0$, we have
\[
b=r+1>r.
\]
Because $b>1>0$, multiplying the inequality $b>1$ by the positive number $b$
gives
\[
b^2>b.
\]
Combining this with $b>r$, we obtain
\[
r<b^2.
\]

Therefore there exist positive rational numbers $a$ and $b$ such that
\[
a^2<r<b^2,
\]
as required.
\end{proof}

\begin{remark*}[Proof structure]
The proof uses an explicit construction rather than an existence argument by
approximation. The lower bracket is chosen as
\[
a=\frac{r}{r+1},
\]
which is both positive and strictly less than $1$ and $r$, so its square is
even smaller. The upper bracket is chosen as
\[
b=r+1,
\]
which is larger than $r$ and larger than $1$, so its square is larger still.
This produces rational squares on both sides of the given positive rational.
\end{remark*}

\begin{dependencies}
\begin{itemize}
  \item Basic field closure of $\mathbb{Q}$ under addition and division
  \item Basic order compatibility in an ordered field
\end{itemize}
\end{dependencies}

\clearpage
